\section{Graph measurements and conditional shrinking}\label{sec:graph}

For a \(C^*\)-algebra \(C\), write \(\widetilde C\) for its
minimal unitization when \(C\) is nonunital, and set
\(\widetilde C=C\) when \(C\) is unital.  We use its canonical
action on \(C^{**}\) in products with support projections.
Given positive contractions \(a,y\in C\), with \(a\) rich,
and a contraction \(c\in\widetilde C\), we seek positive
contractions \(p_1\in\Her(a)\) and \(q_1\in\Her(y)\) with two
properties: each has represented norm one wherever its corresponding
input has norm greater than \(1/2\), and
\(\|e_{p_1}ce_{q_1}\|<1\).  The norm signals will later give
ideal comparison.

The construction has two parts.  A graph of orthogonal copies
first turns possible failure into finitely many explicit equalities.
Once those equalities have been excluded, ideal cutoffs convert
the resulting nonvanishing into one global support estimate.
We prove both parts here while leaving the finite exclusion
problem to the later arguments.  Throughout this section put
\[
 s=\frac14,\qquad \alpha_0=\frac1{32}.
\]
When the integer \(k\ge1\) has been chosen, fix
\(\gamma_g,\sigma_g>0\) with
\begin{equation}\label{eq:graph-constants}
 \gamma_g^2+k\sigma_g^2<1.
\end{equation}
For example, \(\gamma_g=1/2\) and
\(\sigma_g=1/(2\sqrt{k+1})\) satisfy this inequality.  Both the
number of coordinates and these coefficients are fixed before
the measurements below are computed.

\begin{proposition}[Fixed measurements of a graph]\label{prop:graph}
Let \(C\) be a \(C^*\)-algebra, let \(a,y\in C_+\) be
contractions with \(a\) rich, and let \(c\in\widetilde C\) be a
contraction.  For an integer \(k\ge1\) and fixed positive
\(\gamma_g,\sigma_g\) satisfying \eqref{eq:graph-constants},
there are
\[
 X\in\Her((a-1/8)_+)_+,\qquad
 B_1,\ldots,B_k\in\Her((a-1/8)_+)
\]
with the following property.  For any contractions
\[
 S_i\in\Her((X-\alpha_0)_+),\qquad 1\le i\le k,
\]
one can form a rich positive contraction \(p\in\Her(a)\) such
that, for every representation \(\pi\) of \(C\),
\(\|\pi(a)\|>1/2\) implies \(\|\pi(p)\|=1\).
If a representation \(\pi\) of \(C\) satisfies
\[
 \pi\bigl(y(1-c^*p^2c)y\bigr)=0,\qquad \|\pi(y)\|>s,
\]
then, in that same representation,
\begin{equation}\label{eq:graph-failure}
 \|\pi(X)\|>s,\qquad \pi(B_i)=\pi(S_i)\quad(1\le i\le k).
\end{equation}
The elements \(X,B_1,\ldots,B_k\) are fixed before the
\(S_i\) are chosen.
\end{proposition}

\begin{proof}
Let \(e=e_{(a-1/8)_+}\) and \(H=\Her((a-1/8)_+)\).
Since \(a\) is properly infinite,
\(a^{\oplus(k+2)}\precsim a\).  Apply
Lemma~\ref{lem:transport} with cutoff \(1/8\), and split the
resulting row into slots labelled \(0,1,\ldots,k,*\).  They
satisfy
\begin{equation}\label{eq:graph-slots}
 v_i^*v_j=\delta_{ij}e,\qquad v_i v_i^*\le e_a,\qquad
 v_iH\subset C.
\end{equation}
Choose a ramp \(t=t(a)\in H\) that is zero on \([0,1/4]\)
and one on \([1/2,1]\).  It is rich, and its represented norm
is one whenever the represented norm of \(a\) exceeds \(1/2\).

We now compute the measurements, before making any choice of
\(S_i\).  Put \(F=cyc^*\in C_+\).  The block calculation behind
their definition is as follows.  Suppose \(L\in C^{**}\) satisfies
\(L^*L=e\) and has slot coordinates
\(v_0^*L=\gamma_g e\), \(v_i^*L=\sigma_g S_i\), and suppose
\(F=LH_0L^*\) for a positive \(H_0\in eC^{**}e\).  Then
\[
 v_0^*Fv_0=\gamma_g^2H_0,\qquad
 v_i^*Fv_0=\sigma_g\gamma_g S_iH_0.
\]
Vanishing of the defect in a representation will force this form
there, with \(tH_0=H_0\) and with \(t\) acting identically on the
prescribed supports of
\(S_i\).  The first measurement therefore recovers \(H_0\), and
an inverse on its spectral cutoff recovers each \(S_i\) from the
second.  These are conditional identities used to design the
measurements; their verification follows after the graph is built.

Choose a continuous nonnegative function \(h_{\rm inv}\) on
\([0,\infty)\), zero
at zero, equal to \(\lambda^{-1}\) for \(\lambda\ge\alpha_0\),
and bounded by \(\alpha_0^{-1}\).  Define
\begin{equation}\label{eq:graph-measurements}
 X=\gamma_g^{-2}t v_0^*Fv_0t,\qquad
 B_i=(\sigma_g\gamma_g)^{-1}
       t v_i^*Fv_0t\,h_{\rm inv}(X).
\end{equation}
The sandwiches are \((v_it)^*F(v_0t)\), hence belong to \(C\).
They are supported on both sides by \(e\), so they lie in \(H\);
the same is true after multiplication by \(h_{\rm inv}(X)\).
Thus \(X\ge0\) and \(X,B_i\) have the asserted algebra
membership.  The slots, \(t,c,y\) determine them without using
any \(S_i\).

For the chosen \(S_i\), define in the corner \(eC^{**}e\)
\[
 J=\left((1-\gamma_g^2)e
       -\sigma_g^2\sum_{i=1}^kS_i^*S_i\right)^{1/2}.
\]
Then define in \(C^{**}\)
\[
 L=\gamma_g v_0+\sigma_g\sum_{i=1}^k v_iS_i+v_*J,
 \qquad p=LtL^*.
\]
Here \(S_i\in H\), since \(X\in H\).  The radicand is at least
\((1-\gamma_g^2-k\sigma_g^2)e\).  Its square root belongs to
\(H+\mathbb C e\) and multiplies \(H\) into itself.  Hence
\(Lh\in C\) for every \(h\in H\), term by term in
\eqref{eq:graph-slots}.  The slot identities also give
\[
 L^*L=e,\qquad e_aL=L.
\]
It follows, by the membership and support argument of
Lemma~\ref{lem:transport}, that \(h\mapsto LhL^*\) is an
injective \(*\)-homomorphism from \(H\) into \(\Her(a)\).
Indeed, for \(h\ge0\) its value is
\((Lh^{1/2})(Lh^{1/2})^*\in C\), supported under \(e_a\);
and \(L^*L=e\) gives multiplicativity and recovers \(h\) on
conjugating back.  The same identities give preservation of
cutoffs, their Cuntz classes, and represented norms, exactly
as in that lemma.  Thus \(p\) is a rich positive contraction
with the stated representation signal.  If \(e=0\), all the
operators on this corner are zero and these assertions still
hold.

It remains to see what vanishing of the defect forces.  Work
on the essential subspace of a representation in which
\[
 d=y(1-c^*p^2c)y=0,\qquad \|y\|>s,
\]
and omit the representation symbols.  The normal extensions
preserve the preceding support identities.  For every vector
\(\xi\),
\[
 \|y\xi\|=\|pc\,y\xi\|
       \le\|c\,y\xi\|\le\|y\xi\|.
\]
Equality first on \(\ran y\) and then on its closure shows
that \(c\) is isometric on \(e_y\mathcal H\).  For a positive
contraction \(p\), equality \(\|p\zeta\|=\|\zeta\|\) implies
\(p\zeta=\zeta\), by applying \(1-p^2\).  Consequently
\(ce_y\mathcal H\) lies in the \(p=1\) subspace, and the
positive operator \(F=cyc^*\) is supported on that subspace.
Since \(p=LtL^*\) and \(L^*L=e\), the \(p=1\) subspace lies
in the range of \(L\), and its pullback lies in the \(t=1\)
subspace.  With \(H_0=L^*FL\), we obtain
\[
 F=LH_0L^*,\qquad tH_0=H_0=H_0t,\qquad
 \|H_0\|=\|F\|=\|y\|>s.
\]
The last equality uses the isometric action of \(c\) on
\(e_y\mathcal H\).

Now \(v_0^*L=\gamma_g e\), \(v_i^*L=\sigma_g S_i\), and
\(L^*v_0=\gamma_g e\).  Substituting these identities into
\eqref{eq:graph-measurements} yields
\[
 X=tH_0t=H_0,\qquad
 B_i=tS_iH_0t\,h_{\rm inv}(X).
\]
Since \(tX=X\), \(t\) is the identity on
\(\overline{\ran X}=e_X\mathcal H\).  Each represented \(S_i\)
has both supports under \(e_{(X-\alpha_0)_+}\).  Therefore
\(tS_i=S_i\) and \(S_iXh_{\rm inv}(X)=S_i\).  The last display
gives \(B_i=S_i\), while \(\|X\|=\|H_0\|>s\).
This proves \eqref{eq:graph-failure}.
\end{proof}

Thus it is enough to choose the \(S_i\) in their prescribed
spectral corner so that the simultaneous relations
\eqref{eq:graph-failure} cannot hold.  The graph then supplies
the nonvanishing implication assumed in the next lemma.  The
lemma uses the copied element only through its richness and
its representation signal, so it is independent of how the
finite equalities were excluded.

\begin{lemma}[From nonvanishing to shrinking]\label{lem:conditional-shrinking}
Let \(C\) be a \(C^*\)-algebra, let \(a,y\in C_+\) be
contractions with \(a\) rich, and let \(c\in\widetilde C\) be
a contraction.  Suppose \(p\in\Her(a)\) is a rich positive
contraction satisfying
\[
 \|\pi(a)\|>1/2\quad\Longrightarrow\quad\|\pi(p)\|=1
 \qquad\text{for every representation }\pi\text{ of }C.
\]
Put \(d=y(1-c^*p^2c)y\in C_+\), and assume
\begin{equation}\label{eq:common-nonvanishing}
 \pi(d)=0\quad\Longrightarrow\quad\|\pi(y)\|\le\frac14
 \qquad\text{for every representation }\pi\text{ of }C.
\end{equation}
Then there are positive contractions
\[
 p_1\in\Her(a),\qquad q_1\in\Her(y)
\]
such that \(p_1\) is rich,
\begin{equation}\label{eq:shrinking-signals}
 \begin{split}
  \|\pi(a)\|>1/2&\ \Longrightarrow\ \|\pi(p_1)\|=1,\\
  \|\pi(y)\|>1/2&\ \Longrightarrow\ \|\pi(q_1)\|=1
 \end{split}
 \qquad\text{for every representation }\pi\text{ of }C,
\end{equation}
and
\begin{equation}\label{eq:shrinking-support}
 \|e_{p_1}c e_{q_1}\|_{C^{**}}<1.
\end{equation}
The last inequality is a single norm bound in the universal
bidual.
\end{lemma}

\begin{proof}
The contraction assumptions give \(0\le d\le y^2\le1\).
By \eqref{eq:common-nonvanishing}, every irreducible
representation vanishing on \(d\) vanishes on \((y-s)_+\).
The primitive ideal test in Lemma~\ref{lem:ideals-cutoffs}(ii)
therefore gives \((y-s)_+\in\Id(d)\).  Set \(s'=3/8\).
Apply cutoff exhaustion in Lemma~\ref{lem:ideals-cutoffs}(i)
to \((y-s)_+\) with cutoff \(s'-s\).  After decreasing the
resulting parameter if necessary, there is \(0<\delta<1\) with
\begin{equation}\label{eq:shrinking-first-ideal}
 (y-s')_+\in\Id(h_d),\qquad h_d=(d-\delta)_+.
\end{equation}
Define
\[
 z=yh_d^{1/2}\in C,\qquad q=zz^*=y h_d y\in\Her(y).
\]
The order \(d\le y^2\le1\) gives
\begin{equation}\label{eq:shrinking-ideal-order}
 \delta h_d
 \le h_d^{1/2}d h_d^{1/2}
 \le z^*z=h_d^{1/2}y^2h_d^{1/2}
 \le h_d.
\end{equation}
Hence \(z^*z\) and \(h_d\) have the same hereditary algebra,
and \(z^*z\CuEq zz^*=q\).  It follows that
\(\Id(q)=\Id(h_d)\).  A second application of cutoff
exhaustion, to \((y-s')_+\) with cutoff \(1/2-s'\), gives
\(\delta'>0\) such that
\begin{equation}\label{eq:shrinking-second-ideal}
 (y-1/2)_+\in\Id((q-\delta')_+).
\end{equation}
Let \(q_1\) be a ramp of \(q\), zero on \([0,\delta'/2]\)
and one on \([\delta',\infty)\).  It is a positive contraction
in \(\Her(q)\subset\Her(y)\).  If \(\|\pi(y)\|>1/2\),
then \(\pi((y-1/2)_+)\ne0\).  The ideal inclusion
\eqref{eq:shrinking-second-ideal} forces
\(\pi((q-\delta')_+)\ne0\), so \(\|\pi(q)\|>\delta'\)
and \(\|\pi(q_1)\|=1\).

To bound the universal-bidual norm, work in the universal
nondegenerate representation.  For
\(\xi\in e_{h_d}\mathcal H
=\mathbf1_{(\delta,\infty)}(d)\mathcal H\),
\begin{align}
 \|pc\,y\xi\|^2
 &=\|y\xi\|^2-\langle d\xi,\xi\rangle\notag\\
 &\le\|y\xi\|^2-\delta\|\xi\|^2
 \le(1-\delta)\|y\xi\|^2.
 \label{eq:shrinking-norm-loss}
\end{align}
The last inequality uses that \(y\) is a contraction.
Apply this estimate to vectors in \(\ran(h_d^{1/2})\).
The closure of \(\ran(yh_d^{1/2})\) is \(e_q\mathcal H\),
so continuity gives
\[
 \|pc e_q\|\le\sqrt{1-\delta}.
\]
Choose \(\sqrt{1-\delta}<\tau<1\), and take \(p_1\) to be a
ramp of \(p\), zero on \([0,\tau]\) and one at \(1\).
Functional calculus gives \(e_{p_1}\le\tau^{-2}p^2\),
while \(e_{q_1}\le e_q\).  Consequently
\begin{equation}\label{eq:shrinking-final-bound}
 \|e_{p_1}c e_{q_1}\|
 \le\tau^{-1}\|pc e_{q_1}\|
 \le\frac{\sqrt{1-\delta}}{\tau}<1.
\end{equation}
The ramp of the rich \(p\) is rich and remains in \(\Her(a)\).
Spectral mapping at its plateau shows that it has norm one in
every representation where \(p\) has norm one.  This proves
the first signal in \eqref{eq:shrinking-signals}; the second
was proved above.  The proof also covers zero inputs.
\end{proof}
